% ================================================================== Chapitre 1
\chapter{Rappels de première année}

Ce chapitre rassemble, sous forme condensée, les outils de PTSI réutilisés en permanence ensuite : avancement, activité, quotient de réaction et oxydoréduction. Les diagrammes $E$-pH et les titrages rédox, également revus en début de chapitre, ne sont pas repris ici : ils n'interviennent pas dans les raisonnements thermodynamiques qui suivent.

\section{Système physico-chimique et avancement}

\begin{definition}[système physico-chimique, constituant, phase]
\begin{itemize}
  \item Un \textbf{constituant} est une espèce chimique \emph{dans un état physique donné} : \ce{H2O(l)} et \ce{H2O(g)} sont deux constituants différents.
  \item Une \textbf{phase} est une région de l'espace où les grandeurs intensives (masse volumique, concentrations\dots) varient continûment. Un mélange de gaz forme toujours une seule phase ; deux solides distincts forment deux phases.
  \item Un \textbf{système physico-chimique} est un ensemble de constituants répartis en une ou plusieurs phases, susceptible d'être le siège de transformations chimiques ou physiques.
\end{itemize}
\end{definition}

\begin{definition}[équation de réaction, coefficients stœchiométriques algébriques]
On écrit toute équation de réaction sous la forme
\[ 0 = \sum_i \nu_i\,\mathrm{A}_i \]
où $\nu_i$ est le \textbf{coefficient stœchiométrique algébrique} du constituant $\mathrm{A}_i$ : $\nu_i>0$ pour un produit, $\nu_i<0$ pour un réactif, $\nu_i = 0$ pour une espèce spectatrice.
\end{definition}

\begin{exemple}
Pour \ce{N2(g) + 3 H2(g) = 2 NH3(g)} : $\nu_{\ce{N2}} = -1$, $\nu_{\ce{H2}} = -3$, $\nu_{\ce{NH3}} = +2$.
\end{exemple}

\begin{definition}[avancement]
L'\textbf{avancement} $\xi$ (en mol) est défini par
\[ n_i(\xi) = n_{i,0} + \nu_i\,\xi \qquad\Longleftrightarrow\qquad \dd n_i = \nu_i\,\dd\xi \quad\text{pour tout }i. \]
Il est positif si la réaction avance dans le sens direct ($\rightarrow$), négatif dans le sens indirect ($\leftarrow$). L'\textbf{avancement maximal} $\xi_{\max}$ correspond à l'épuisement du réactif limitant.
\end{definition}

\section{Activité, quotient de réaction et constante d'équilibre}

\begin{definition}[activité]
L'\textbf{activité} $a_i$ d'un constituant est une grandeur sans dimension qui mesure son « poids » effectif dans le mélange. Dans le cadre du programme (cas idéaux) :

\renewcommand{\arraystretch}{1.3}
\begin{center}
\begin{tabular}{ll}
\toprule
\textbf{Constituant} & \textbf{Activité} \\ \midrule
gaz parfait dans un mélange & $a_i = \dfrac{P_i}{P^\circ} = x_i\dfrac{P}{P^\circ}$ \quad ($P_i$ pression partielle) \\[0.6em]
soluté en solution diluée & $a_i = \dfrac{c_i}{c^\circ}$ \\[0.6em]
solvant (eau) & $a = 1$ \\
solide ou liquide \emph{pur} (seul dans sa phase) & $a = 1$ \\
\bottomrule
\end{tabular}
\end{center}
\end{definition}

\begin{definition}[quotient de réaction]
Pour la réaction $0 = \sum_i \nu_i\mathrm{A}_i$, le \textbf{quotient de réaction} est
\[ Q_r = \prod_i a_i^{\,\nu_i} \qquad\text{(produits au numérateur, réactifs au dénominateur).} \]
Il dépend de l'état du système : il varie au cours de la transformation.
\end{definition}

\begin{rappel}[loi d'action des masses et critère d'évolution]
\begin{itemize}
  \item À l'équilibre chimique, $Q_{r,\text{éq}} = K^\circ(T)$ : la \textbf{constante thermodynamique d'équilibre} $K^\circ$ ne dépend que de la température.
  \item Si $Q_r < K^\circ$, le système évolue dans le sens direct ; si $Q_r > K^\circ$, dans le sens indirect.
\end{itemize}
En PTSI ce critère était \emph{admis}. Le chapitre~\ref{ch:evolution} le \emph{démontre} à partir du deuxième principe, et les chapitres~\ref{ch:grandeurs} et~\ref{ch:evolution} montrent comment \emph{calculer} $K^\circ$ (qui était toujours donnée en première année).
\end{rappel}

\section{Oxydoréduction}

\subsection{Couples et nombres d'oxydation}

\begin{definition}[oxydant, réducteur, demi-équation]
Un \textbf{oxydant} est une espèce capable de capter des électrons, un \textbf{réducteur} une espèce capable d'en céder. Ils forment un \textbf{couple} Ox/Red relié par la demi-équation électronique
\[ \alpha\,\mathrm{Ox} + \gamma\,e^- \;=\; \beta\,\mathrm{Red} \]
(équilibrée si besoin avec \ce{H+} et \ce{H2O}). Une \textbf{réaction d'oxydoréduction} échange $\gamma$ électrons entre l'oxydant d'un couple et le réducteur d'un autre.
\end{definition}

\begin{definition}[nombre d'oxydation]
Le \textbf{nombre d'oxydation} (n.o.) d'un élément dans un édifice est la charge fictive qu'il porterait si l'on attribuait chaque doublet de liaison à l'atome le plus électronégatif. On le note en chiffres romains. Règles pratiques :
\begin{itemize}
  \item la somme des n.o. des atomes est égale à la charge de l'édifice ;
  \item O vaut en général $-\mathrm{II}$ (sauf peroxydes : $-\mathrm{I}$) ; H vaut en général $+\mathrm{I}$ ;
  \item un corps simple (\ce{O2}, \ce{Fe}, \ce{H2}\dots) a tous ses n.o. nuls.
\end{itemize}
Une \emph{oxydation} fait \emph{croître} le n.o. ; une \emph{réduction} le fait \emph{décroître}.
\end{definition}

\begin{exemple}[calculs de n.o.]
\begin{itemize}
  \item \ce{ClO4^-} : $x + 4\times(-2) = -1$ donc $x = +\mathrm{VII}$ pour le chlore (on retrouve le résultat en attribuant tous les doublets Cl--O à l'oxygène dans la formule de Lewis).
  \item \ce{H2O2} : $2\times(+1) + 2x = 0$ donc O est au n.o. $-\mathrm{I}$ (pont peroxyde \ce{O-O}).
  \item \ce{O2^-} (ion superoxyde) : $2x = -1$ donc $x = -\tfrac12$. Un n.o. non entier est une \emph{moyenne} : en réalité, les deux atomes ne sont pas dans le même environnement électronique.
\end{itemize}
\end{exemple}

\subsection{Potentiel d'électrode et formule de Nernst}

\begin{definition}[électrode, potentiel d'électrode]
Une \textbf{électrode} (ou demi-pile) est l'ensemble formé d'un conducteur électronique (métal, graphite\dots) au contact d'un électrolyte contenant les deux membres d'un couple. Son \textbf{potentiel} $E$ est la tension entre ce conducteur et une électrode de référence : l'\textbf{électrode standard à hydrogène} (ESH), dont le potentiel est nul par convention à toute température.
\end{definition}

\begin{resultat}{Formule de Nernst}
Pour le couple $\alpha\,\mathrm{Ox} + \gamma e^- + \dots = \beta\,\mathrm{Red} + \dots$ à l'équilibre électrochimique :
\[ E = E^\circ(\mathrm{Ox/Red}) + \frac{RT}{\gamma F}\ln\frac{a_{\mathrm{Ox}}^{\,\alpha}\,(\dots)}{a_{\mathrm{Red}}^{\,\beta}\,(\dots)} \;\underset{\SI{298}{K}}{\approx}\; E^\circ + \frac{0{,}06}{\gamma}\log\frac{a_{\mathrm{Ox}}^{\,\alpha}\,(\dots)}{a_{\mathrm{Red}}^{\,\beta}\,(\dots)} \]
où $E^\circ$ est le \textbf{potentiel standard} du couple (il ne dépend que de $T$) et où les « $(\dots)$ » désignent les activités des autres espèces de la demi-équation (\ce{H+}, \ce{Cl-}\dots), chacune à la puissance de son coefficient.
\end{resultat}

\begin{exemple}[quelques électrodes classiques]
\begin{itemize}
  \item \textbf{Métal dans une solution de ses ions} : \ce{Cu^{2+} + 2e^- = Cu(s)}, donc $E = 0{,}34 + 0{,}03\log\dfrac{[\ce{Cu^{2+}}]}{c^\circ}$.
  \item \textbf{Électrode à gaz} : \ce{2H+ + 2e^- = H2(g)} sur platine, donc $E = 0 + 0{,}03\log\dfrac{([\ce{H+}]/c^\circ)^2}{P_{\ce{H2}}/P^\circ}$. Dans les conditions standard ($\mathrm{pH} = 0$, $P_{\ce{H2}} = \SI{1}{bar}$), $E = \SI{0}{V}$ : c'est l'ESH.
  \item \textbf{Électrode de deuxième espèce} (métal, sel peu soluble de ce métal, anion commun) : l'électrode au calomel \ce{Hg2Cl2(s)/Hg(l)} dans KCl, de potentiel $E = 0{,}27 - 0{,}06\log\dfrac{[\ce{Cl-}]}{c^\circ}$. Dans KCl saturé, c'est l'électrode au calomel saturée (ECS), $E_{\text{ECS}} \approx \SI{0,24}{V}$. Le calcul du « $0{,}27$ » est fait au chapitre~\ref{ch:redox}.
\end{itemize}
\end{exemple}

\subsection{Prévision des réactions : règle du gamma}

\begin{figure}[H]
\centering
\begin{tikzpicture}[>=Stealth]
  \draw[->,line width=0.9pt] (0,-0.3) -- (0,4.5) node[above] {$E^\circ$};
  \draw (-0.12,3.2) -- (0.12,3.2); \draw (-0.12,0.8) -- (0.12,0.8);
  \node[left] at (-0.25,3.2) {$\mathrm{Ox}_1$}; \node[right] at (0.25,3.2) {$\mathrm{Red}_1$};
  \node[left] at (-0.25,0.8) {$\mathrm{Ox}_2$}; \node[right] at (0.25,0.8) {$\mathrm{Red}_2$};
  \node[left,font=\small,color=black!60] at (-1.4,3.2) {$E_1^\circ$};
  \node[left,font=\small,color=black!60] at (-1.4,0.8) {$E_2^\circ$};
  \draw[c3,line width=1.2pt,->] (-0.75,3.45) to[out=30,in=150] node[above,font=\scriptsize,fill=white,inner sep=1pt] {réduction} (0.85,3.45);
  \draw[c4,line width=1.2pt,->] (0.85,0.55) to[out=-150,in=-30] node[below,font=\scriptsize,fill=white,inner sep=1pt] {oxydation} (-0.75,0.55);
  \draw[navy,line width=1pt,dashed] (-0.7,2.95) -- (0.75,1.05);
  \node[c3,font=\small,align=left,anchor=west] at (2.0,2.0) {réaction spontanée (sens direct, $K^\circ>1$) :\\[0.2em]$\mathrm{Ox}_1 + \mathrm{Red}_2 \longrightarrow \mathrm{Red}_1 + \mathrm{Ox}_2$};
\end{tikzpicture}
\caption{Règle du gamma : l'oxydant le plus fort (en haut à gauche) réagit sur le réducteur le plus fort (en bas à droite).}
\label{fig:gamma}
\end{figure}

\begin{resultat}{Constante d'équilibre d'une réaction d'oxydoréduction (PTSI)}
Pour $a_1\mathrm{Ox}_1 + b_2\mathrm{Red}_2 = b_1\mathrm{Red}_1 + a_2\mathrm{Ox}_2$ échangeant $\gamma$ électrons :
\[ \ln K^\circ = \frac{\gamma F}{RT}\left(E_1^\circ - E_2^\circ\right) \qquad\Longrightarrow\qquad \boxed{\log K^\circ \underset{\SI{298}{K}}{=} \frac{\gamma\,(E_1^\circ - E_2^\circ)}{0{,}06}} \]
\end{resultat}

\begin{demo}
À l'équilibre, la solution n'a qu'un seul potentiel : $E_1 = E_2$. En écrivant Nernst pour les deux couples (avec le même nombre $\gamma$ d'électrons, après multiplication des demi-équations) :
\[ E_1^\circ + \frac{RT}{\gamma F}\ln\frac{a_{\mathrm{Ox}_1}^{a_1}}{a_{\mathrm{Red}_1}^{b_1}} = E_2^\circ + \frac{RT}{\gamma F}\ln\frac{a_{\mathrm{Ox}_2}^{a_2}}{a_{\mathrm{Red}_2}^{b_2}} \;\Longrightarrow\; E_1^\circ - E_2^\circ = \frac{RT}{\gamma F}\ln\underbrace{\frac{a_{\mathrm{Red}_1}^{b_1}\,a_{\mathrm{Ox}_2}^{a_2}}{a_{\mathrm{Ox}_1}^{a_1}\,a_{\mathrm{Red}_2}^{b_2}}}_{Q_{r,\text{éq}} = K^\circ} \]
\end{demo}

Le chapitre~\ref{ch:redox} retrouve ce résultat par une voie purement thermodynamique, en reliant $\Delta_rG$ aux potentiels.
